\chapter{Quantum Fields in Curved Spacetime}
\label{ch:qft-curved}

\section{The General Action and Equation of Motion}

We consider a free scalar field on a curved spacetime background,
described by the action
\begin{equation}
    S = \int d^dx\sqrt{-g}\left[
    -\frac{1}{2}g^{\mu\nu}\partial_\mu\phi\,\partial_\nu\phi
    - \frac{1}{2}(m^2 - \xi R)\phi^2\right],
\end{equation}
where $R$ is the Ricci scalar and $\xi$ is a non-minimal coupling to
the curvature. The special value
\begin{equation}
    \xi = \frac{d-2}{4(d-1)}
\end{equation}
is the conformal coupling, for which the massless action acquires an
additional symmetry: it becomes invariant under a simultaneous scaling
of the metric and the field. The equation of motion is
\begin{equation}
    \left(\square - m^2 + \xi R\right)\phi = 0,
    \qquad \square = g^{\mu\nu}\nabla_\mu\nabla_\nu.
\end{equation}
In flat spacetime this reduces to the Klein-Gordon equation, which we
review first before treating the curved case.

\section{Quantization in Flat Spacetime}

In 4D Minkowski spacetime the field is expanded in plane wave modes,
\begin{equation}
    \phi(t,\mathbf{x}) = \int d^3k
    \left(u_k(t,\mathbf{x})\,a_k
    + u_k^*(t,\mathbf{x})\,a_k^\dagger\right),
\end{equation}
where the positive and negative frequency modes are
\begin{align}
    u_k &= \frac{1}{(2\pi)^{3/2}\sqrt{2\omega}}
    e^{-i\omega t + i\mathbf{k}\cdot\mathbf{x}}, \\
    u_k^* &= \frac{1}{(2\pi)^{3/2}\sqrt{2\omega}}
    e^{+i\omega t - i\mathbf{k}\cdot\mathbf{x}},
\end{align}
with dispersion relation $\omega = +\sqrt{\mathbf{k}^2+m^2}$. The
canonical conjugate momentum is
\begin{equation}
    \Pi(t,\mathbf{x})
    = \frac{\partial\mathcal{L}}{\partial\dot\phi}
    = \dot\phi(t,\mathbf{x}),
\end{equation}
and the equal-time commutation relations are
\begin{align}
    \left[\Pi(t,\mathbf{x}),\phi(t,\mathbf{x}')\right]
    &= -i\,\delta(\mathbf{x}-\mathbf{x}'), \\
    \left[\Pi(t,\mathbf{x}),\Pi(t,\mathbf{x}')\right]
    &= \left[\phi(t,\mathbf{x}),\phi(t,\mathbf{x}')\right] = 0.
\end{align}
The vacuum $|0\rangle$ is defined by $a_k|0\rangle = 0$ for all $k$,
and $n$-particle states are built by acting with creation operators
$a_k^\dagger$.

The modes are organized by the Klein-Gordon inner product, defined for
two solutions $\varphi_1$, $\varphi_2$ of the field equation as
\begin{equation}
    \langle\varphi_1,\varphi_2\rangle
    \equiv -i\int_{\text{fixed }t}d^3x
    \left(\varphi_1\,\partial_t\varphi_2^*
    - \partial_t\varphi_1\cdot\varphi_2^*\right).
\end{equation}
This product is conserved in time, as follows from the Klein-Gordon
equation. It is not positive definite on the space of all solutions,
but it is positive definite on the space of positive-frequency
solutions and negative definite on negative-frequency solutions, with
the two spaces orthogonal to each other. The orthonormality relations
for the plane wave modes are
\begin{align}
    \langle u_k, u_{k'}\rangle &= \delta(\mathbf{k}-\mathbf{k}'), \\
    \langle u_k^*, u_{k'}^*\rangle &= -\delta(\mathbf{k}-\mathbf{k}'), \\
    \langle u_k, u_{k'}^*\rangle &= 0.
\end{align}
We verify the first relation explicitly. Computing the time derivatives
$\partial_t u_{k'}^* = +i\omega' u_{k'}^*$ and
$\partial_t u_k = -i\omega u_k$,
\begin{align}
    \langle u_k, u_{k'}\rangle
    &= -i\int d^3x\,u_k(+i\omega')u_{k'}^*
    - u_{k'}^*(-i\omega)u_k \notag\\
    &= (\omega+\omega')\int d^3x\,u_k\,u_{k'}^* \notag\\
    &= \frac{\omega+\omega'}{(2\pi)^3\cdot 2\sqrt{\omega\omega'}}
    \int d^3x\,e^{i(\mathbf{k}-\mathbf{k}')\cdot\mathbf{x}} \notag\\
    &= \frac{\omega+\omega'}{2\sqrt{\omega\omega'}}
    \,\delta(\mathbf{k}-\mathbf{k}').
\end{align}
Since $\delta(\mathbf{k}-\mathbf{k}')\neq 0$ only when
$\mathbf{k}=\mathbf{k}'$, i.e.\ $\omega=\omega'$, the prefactor
evaluates to $2\omega/(2\omega) = 1$, confirming the first relation.
The second follows from antisymmetry of the inner product, and the
third from the fact that
$\langle u_k,u_{k'}^*\rangle \propto (\omega-\omega')
\delta(\mathbf{k}+\mathbf{k}')$, which vanishes when
$\mathbf{k}'=-\mathbf{k}$ forces $\omega=\omega'$. These relations
allow the creation and annihilation operators to be extracted directly
from the field,
\begin{equation}
    a_k = \langle\phi, u_k\rangle, \qquad
    a_k^\dagger = -\langle\phi, u_k^*\rangle,
\end{equation}
with commutation relations $[a_k, a_{k'}^\dagger]
= \langle u_k, u_{k'}\rangle = \delta(\mathbf{k}-\mathbf{k}')$.

\section{Quantization in Curved Spacetime}

In a curved spacetime we choose a spacelike foliation: a family of
spacelike hypersurfaces $\Sigma$ parameterized by a time function
$t(x^\mu)$.

% [FIGURE: spacelike foliation of a curved spacetime]

The unit future-pointing normal to $\Sigma$ is
\begin{equation}
    n_\mu = \frac{\partial_\mu t}{|\partial_\mu t|},
    \qquad n_\mu n_\nu g^{\mu\nu} = -1,
\end{equation}
and the induced metric on $\Sigma$ is
\begin{equation}
    (g_\Sigma)_{\mu\nu} = g_{\mu\nu} + n_\mu n_\nu,
\end{equation}
which satisfies $(g_\Sigma)_{\mu\nu}n^\mu n^\nu = 0$, confirming that
it projects vectors onto the surface. The Klein-Gordon inner product
generalizes to
\begin{equation}
    \langle\varphi_1,\varphi_2\rangle
    = -i\int_\Sigma d^{d-1}\sigma\sqrt{g_\Sigma}\,n^\mu
    \left(\varphi_1\,\partial_\mu\varphi_2^*
    - \partial_\mu\varphi_1\cdot\varphi_2^*\right).
\end{equation}
This inner product is independent of the choice of spacelike surface
$\Sigma$, as follows from the Gauss divergence theorem. If $\Sigma_1$
and $\Sigma_2$ are two spacelike slices bounding a four-volume
$\Omega$, then
\begin{align}
    \langle\varphi_1,\varphi_2\rangle_{\Sigma_1}
    - \langle\varphi_1,\varphi_2\rangle_{\Sigma_2}
    &= \int_\Omega d^4x\sqrt{g}\,
    \nabla^\mu\!\left(\varphi_1^*\nabla_\mu\varphi_2
    - \nabla_\mu\varphi_1^*\cdot\varphi_2\right) = 0,
\end{align}
where the last equality follows because both $\varphi_1$ and
$\varphi_2$ satisfy the Klein-Gordon equation. The inner product
therefore has the same properties as in flat space, provided the
contribution from spatial infinity is negligible.

The field is expanded in a complete set of modes $\{u_i, u_i^*\}$
satisfying the same orthonormality relations as before,
\begin{equation}
    \phi(x) = \sum_i\left(u_i(x)\,a_i + u_i^*(x)\,a_i^\dagger\right),
\end{equation}
\begin{equation}
    \langle u_i,u_j\rangle = \delta_{ij}, \qquad
    \langle u_i^*,u_j^*\rangle = -\delta_{ij}, \qquad
    \langle u_i,u_j^*\rangle = 0.
\end{equation}
These are no longer plane waves; their form depends on the spacetime
geometry. The canonical momentum is
\begin{equation}
    \Pi = \frac{\partial\mathcal{L}}{\partial(\partial_t\phi)}
    = \sqrt{g_\Sigma}\,n^\mu\partial_\mu\phi,
\end{equation}
and the equal-time commutation relations take the same form as in flat
space. The creation and annihilation operators are extracted as
\begin{align}
    a_i &= \langle\phi,u_i\rangle
    = -i\int d\sigma\left(\sqrt{g_\Sigma}\,n^\mu\partial_\mu u_i^*
    \cdot\phi - u_i^*\Pi\right), \\
    a_i^\dagger &= -\langle\phi,u_i^*\rangle
    = i\int d\sigma\left(\sqrt{g_\Sigma}\,n^\mu\partial_\mu u_i
    \cdot\phi - u_i\Pi\right),
\end{align}
satisfying $[a_i,a_j^\dagger] = \delta_{ij}$. The vacuum $|0\rangle$
is defined by $a_i|0\rangle = 0$.

\section{Ambiguity in the Notion of Particles and Bogoliubov
Transformations}

The mode decomposition above is not unique. Consider a second complete
set of modes $\{\tilde u_j, \tilde u_j^*\}$ with the same
orthonormality properties, giving a second quantization with operators
$\hat a_j$ and vacuum $|\tilde 0\rangle$ defined by
$\hat a_j|\tilde 0\rangle = 0$. Since both sets are complete, we can
expand one in terms of the other,
\begin{align}
    \tilde u_j(x) &= \sum_i
    \left(\alpha_{ji}\,u_i(x) + \beta_{ji}\,u_i^*(x)\right), \\
    \tilde u_j^*(x) &= \sum_i
    \left(\alpha_{ji}^*\,u_i^* + \beta_{ji}^*\,u_i\right),
\end{align}
where the Bogoliubov coefficients are
\begin{equation}
    \alpha_{ji} = \langle\tilde u_j, u_i\rangle, \qquad
    \beta_{ji} = -\langle\tilde u_j, u_i^*\rangle.
\end{equation}
The orthonormality of both sets implies the relations
\begin{align}
    \sum_j\left(\alpha_{ji}^*\alpha_{jk} - \beta_{ji}\beta_{jk}^*\right)
    &= \delta_{ik}, \\
    \sum_j\left(\alpha_{ji}^*\beta_{jk} - \beta_{ji}\alpha_{jk}^*\right)
    &= 0,
\end{align}
and their transposes. These can be used to invert the expansion,
\begin{equation}
    u_i = \sum_j\left(\alpha_{ji}\,\tilde u_j
    - \beta_{ji}\,\tilde u_j^*\right).
\end{equation}
The two sets of operators are related by the Bogoliubov transformation,
\begin{align}
    a_i &= \sum_j\left(\alpha_{ji}\,\hat a_j
    + \beta_{ji}^*\,\hat a_j^\dagger\right), \\
    \hat a_j &= \sum_i\left(\alpha_{ji}^*\,a_i
    - \beta_{ji}^*\,a_i^\dagger\right).
\end{align}
The two quantizations are inequivalent precisely when the $\beta$
coefficients are nonzero. To see what this means physically, compute
the expectation value of the number operator $\hat N_j
= \hat a_j^\dagger\hat a_j$ in the vacuum $|0\rangle$ of the first
quantization,
\begin{align}
    \langle 0|\hat a_j^\dagger\hat a_j|0\rangle
    &= \langle 0|\sum_i\left(\alpha_{ji}\,a_i^\dagger
    - \beta_{ji}\,a_i\right)
    \sum_k\left(\alpha_{jk}^*\,a_k
    - \beta_{jk}^*\,a_k^\dagger\right)|0\rangle \notag\\
    &= \sum_{i,k}\beta_{ji}\beta_{jk}^*
    \underbrace{\langle 0|a_i\,a_k^\dagger|0\rangle}_{\delta_{ik}}
    \notag\\
    &= \sum_i|\beta_{ji}|^2.
\end{align}
The state $|0\rangle$, which contains no particles of type $a$,
contains $\sum_i|\beta_{ji}|^2$ particles of type $\hat a_j$. This is
the key result: different choices of modes, corresponding to different
observers or different regions of spacetime, lead to genuinely
inequivalent notions of particle and vacuum. The particle content of a
state is observer-dependent in curved spacetime, in sharp contrast with
the situation in flat spacetime where all inertial observers agree.

In summary, the field equation $(\square - m^2 + \xi R)\phi = 0$ in
curved spacetime admits multiple inequivalent quantizations related by
Bogoliubov transformations,
\begin{equation}
    \phi = \sum_i\left(u_i\,a_i + u_i^*\,a_i^\dagger\right)
    = \sum_j\left(\tilde u_j\,\hat a_j
    + \tilde u_j^*\,\hat a_j^\dagger\right),
\end{equation}
and the vacuum of one observer contains particles as seen by another
whenever $\beta_{ji}\neq 0$. The physical consequences of this
ambiguity are explored in the following chapters.
